% !TeX root = ../icml2026.tex
\section{Conclusion and future work}\label{sec:conclusion}

This paper establishes the theoretical bounds for the minimal embeddable dimensions and studies an additional centroid setting. It revealed a positive fact: we can use a small number of  dimensions to support an extensive retrieval system, provided we consider only queries with small cardinality. Meanwhile, we can clearly conclude that the effectiveness of embedding-based retrieval systems is limited by \textbf{learnability} (the way we build the specific retrieval system) and debunk the claim from the \textbf{approximability} that the theoretical limitation lies in the fundamental property of the vector space that prevents the subset membership structure from being approximated.

Future work should include more discussion of advanced embedding spaces, such as hyperbolic and Wasserstein spaces; a more sophisticated structure for queries and their answers, including situations where the cardinality of answer sets follows a power-law distribution; and a more realistic setting of fixed-point float numbers, such as float point 32, or even float point 4.
